% ID: SEL-D0011-BD-E07
% ESTADO: validada. S012; tres pasadas del mismo autor, sin segundo revisor.
\textbf{a) Convenciones de muestreo y probabilidad.}\par\nopagebreak[4]
El enunciado no especifica reemplazamiento ni orden. Los criterios puntúan el número de muestras y la probabilidad, sin fijar una convención. Explicamos las alternativas habituales.
\emph{Sin reemplazamiento y sin orden}, el muestreo aleatorio simple da
\[
\boxed{\binom92=36\text{ muestras equiprobables}.}
\]
La media es $5$ cuando la suma es $10$: las muestras favorables son $\{1,9\},\{2,8\},\{3,7\},\{4,6\}$. Así,
\[
\boxed{P(\overline X=5)=\frac4{36}=\frac19.}
\]
\emph{Sin reemplazamiento y registrando el orden}, hay $9\cdot8=72$ pares equiprobables y $8$ favorables: la probabilidad sigue siendo $1/9$.
\emph{Con reemplazamiento y extracciones independientes}, hay $9^2=81$ pares ordenados equiprobables. A los ocho favorables anteriores se añade $(5,5)$: $9/81=1/9$. Si esos pares se escriben sin orden hay $45$ tipos, pero no son equiprobables: los distintos pesan $2/81$ y los repetidos $1/81$. No se puede dividir el número de tipos favorables entre $45$ sin ponderar. El número de muestras depende de la convención; la probabilidad pedida coincide en estos modelos.

\textbf{b.1) Intervalo para la proporción.}\par\nopagebreak[4]
Con $n=10000$ y $500$ positivos, $\widehat p=0.05$. Para confianza 97\,\%, $z_{0.985}\approx2.170090378$. El intervalo normal aproximado es
\[
0.05\pm z_{0.985}\sqrt{\frac{0.05\cdot0.95}{10000}},
\]
\[
\boxed{IC_{\text{97\%}}\approx(0.045270398,\ 0.054729602).}
\]
Los recuentos $500$ y $9500$ permiten la aproximación normal habitual, bajo muestreo aleatorio representativo e independencia en el modelo. La confianza expresa la cobertura del procedimiento.

\textbf{b.2) Compatibilidad con el 6\,\%.}\par\nopagebreak[4]
$0.06$ está por encima del extremo superior, aproximadamente $0.05473$. Por tanto, \textbf{no es compatible con el intervalo al nivel indicado}, bajo las hipótesis del método. No se presenta esa exclusión estadística como una imposibilidad absoluta.

\textbf{b.3) Muestra mayor con igual proporción.}\par\nopagebreak[4]
El centro y el cuantil se mantienen; el error es proporcional a $1/\sqrt n$. Si $n'>n$, entonces $E'<E$. El nuevo intervalo \textbf{está estrictamente contenido en el anterior}. La conclusión depende de mantener exactamente la misma proporción y confianza, tal como establece la pregunta.
